\documentclass[11pt]{article}
\usepackage{times}
\usepackage[margin=0.9in]{geometry}
\usepackage{parskip}
\usepackage{hyperref}
\pagestyle{empty}

\begin{document}
\enlargethispage{1.5\baselineskip}

\begin{flushleft}

Yichao Jin, Ph.D.\\
School of Economic, Political and Policy Sciences\\
University of Texas at Dallas\\
yichao.jin@utdallas.edu\\

\vspace{12pt}

August 31, 2026\\

\vspace{12pt}

Editorial Office\\
PharmacoEconomics\\
Springer\\

\vspace{12pt}

Dear Editor,\\

\vspace{12pt}

I am pleased to submit our manuscript, ``Cost-Effectiveness of Intismeran Autogene Plus Pembrolizumab Versus Pembrolizumab Alone as Adjuvant Therapy for Resected Stage IIIB--IV Melanoma: An Updated Early Economic Evaluation,'' for consideration in the \textit{PharmacoEconomics} Gene and Cell Therapies collection.

Individualized neoantigen therapies (INTs) represent a new class of cancer treatment with a cost structure fundamentally different from conventional oncology drugs---each dose is custom-manufactured for a single patient. Our analysis estimates the cost-effectiveness of intismeran (mRNA-4157/V940) plus pembrolizumab, the first INT to demonstrate Phase 3 efficacy in the INTerpath-001 trial. Using a four-state partitioned survival model calibrated to the 5-year KEYNOTE-942 data, we find that the combination is cost-effective at an assumed acquisition price of \$200,000 per course, with an ICER of \textbf{\$42,499/QALY} and \textbf{94.8\%} probability of cost-effectiveness at \$100,000/QALY. A deterministic threshold price analysis shows the ICER remains below \$100,000/QALY at prices up to approximately \$465,669 per course.

We believe this work is well suited to the Gene and Cell Therapies collection for three reasons. First, it addresses the economic evaluation of a novel mRNA-based personalized therapy---a category of treatment that conventional CEA frameworks have not yet fully accommodated. Second, it provides a publicly available model and codebase (\url{https://github.com/yichao2022/cea-intismeran}), contributing to the methodological infrastructure needed for transparent assessment of gene and cell therapies. Third, the threshold price analysis offers practical insight for payers evaluating value-based pricing of individualized therapies.

The manuscript has not been published elsewhere and is not under consideration by another journal. The author declares no competing interests. All data used in this analysis are derived from published sources cited in the references. The model code is freely available at the GitHub repository above.

Thank you for your consideration.

\vspace{20pt}

Sincerely,

\vspace{20pt}

Yichao Jin, Ph.D.\\
\
School of Economic, Political and Policy Sciences\\
University of Texas at Dallas\\
yichao.jin@utdallas.edu\\
https://yichao2022.github.io\\

\end{flushleft}

\end{document}
